\section{Recovery of the Standard Cosmological Model}

\subsection{The Low-Energy Limit of the Hyper-Electron Hamiltonian}

To establish the viability of Excited-State Cosmology, we must rigorously indicate that the high-energy, hyper-dimensional quantum formalism reduces to the classical Einstein-Hilbert action of General Relativity in the low-energy macroscopic limit. This procedure relies on the quantum-mechanical correspondence principle and the path integral formulation of the hyper-electron effective action.

The total quantum action of the hyper-electron system is defined over the higher-dimensional bulk space $\mathcal{M}_{\text{bulk}}$. The state of the system is governed by the hyper-Dirac action, which includes the kinetic term, the central Coulombic hyper-potential $V(r)$, and an emergent gauge field coupling describing the internal 3-manifold topology. The generating functional $Z$ is written as:
\begin{equation}
Z = \int \mathcal{D}[\Psi] \mathcal{D}[\bar{\Psi}] \mathcal{D}[h_{\mu\nu}] \exp\left( \frac{i}{\hbar} S_{\text{hyper}}[\Psi, \bar{\Psi}, h_{\mu\nu}] \right)
\end{equation}
where $\Psi$ is the hyper-fermion field and $h_{\mu\nu}$ represents the induced metric tensor on the 3-brane mapping the internal observable universe.

In the low-energy limit ($\hbar \to 0$, $r \gg a_0$), the high-frequency quantum modes of the hyper-electron state are heavily suppressed, and the saddle-point approximation dominates the path integral. We perform a Kaluza-Klein reduction of the bulk geometry by integrating out the extra dimensions orthogonal to the internal 3-brane. The effective action on the 4-dimensional spacetime $\mathcal{M}_4$ (one temporal and three spatial dimensions) is extracted from the trace anomaly of the hyper-stress-energy tensor:
\begin{equation}
S_{\text{eff}} = \int_{\mathcal{M}_4} d^4x \sqrt{-g} \langle \hat{T}^{\mu}_{\mu} \rangle_{\text{bulk}}
\end{equation}

Using the Schwinger-DeWitt proper-time expansion, the expectation value of the stress-energy tensor yields a leading-order term proportional to the intrinsic Ricci scalar $R$ of the induced metric $g_{\mu\nu}$, and a constant term arising from the ground-state vacuum energy of the hyper-electron zero-point oscillations. The expansion gives:
\begin{equation}
S_{\text{eff}} \approx \int d^4x \sqrt{-g} \left[ \frac{1}{16\pi G_{\text{eff}}} R - \Lambda_{\text{eff}} + \mathcal{L}_{\text{matter}} \right]
\end{equation}
Here, the effective Newton's constant $G_{\text{eff}}$ is inversely proportional to the hyper-mass $\mu$ projected onto the brane, and the cosmological constant $\Lambda_{\text{eff}}$ is precisely the energy eigenvalue difference corresponding to the spontaneous emission decay rate discussed in previous sections. This result formally constitutes the classical Einstein-Hilbert action with a cosmological constant.

\subsection{Recovery of Local General Relativity and $\Lambda$CDM}

To guarantee that local observations match the predictions of standard General Relativity, we examine the behavior of the effective action under local diffeomorphism invariance. Because the dimensional reduction procedure preserves the $GL(4, \mathbb{R})$ symmetry of the induced metric on the internal manifold, the classical equations of motion derived from varying $S_{\text{eff}}$ with respect to $g^{\mu\nu}$ are the standard Einstein field equations:
\begin{equation}
R_{\mu\nu} - \frac{1}{2} R g_{\mu\nu} + \Lambda_{\text{eff}} g_{\mu\nu} = 8\pi G_{\text{eff}} T_{\mu\nu}
\end{equation}

At local scales (e.g., stellar and galactic environments), the internal energy density of baryonic and dark matter dominates the homogeneous background energy $\Lambda_{\text{eff}}$. In the weak-field limit, the metric perturbation $g_{\mu\nu} = \eta_{\mu\nu} + h_{\mu\nu}$ leads directly to the Poisson equation for Newtonian gravity, $\nabla^2 \Phi = 4\pi G_{\text{eff}} \rho$, confirming the complete recovery of classical gravitational physics.

Furthermore, on cosmological scales, assuming spatial homogeneity and isotropy as dictated by the $2p$ azimuthal symmetry in the long-wavelength limit, the metric assumes the Friedmann-Lema\^itre-Robertson-Walker (FLRW) form. Inserting the FLRW metric into the derived field equations yields the standard Friedmann equations:
\begin{equation}
H^2 = \left(\frac{\dot{a}}{a}\right)^2 = \frac{8\pi G_{\text{eff}}}{3} \rho + \frac{\Lambda_{\text{eff}}}{3} - \frac{k}{a^2}
\end{equation}
With the vacuum energy explicitly sourced by the hyper-electron excited state stability (providing $\Lambda_{\text{eff}}$) and the matter density $\rho$ mapping to internal field excitations, this formulation exactly reproduces the phenomenological evolution of the $\Lambda$CDM model. The underlying quantum mechanical structure of the universe is thus safely concealed at low energies, validating the macroscopic equivalence between Excited-State Cosmology and standard cosmological physics.
